\documentclass[10pt,a4paper]{amsart}
\usepackage[margin=19mm]{geometry}
\usepackage{amsmath,amssymb,mathtools}
\usepackage[scr=rsfs,cal=euler]{mathalfa}
\usepackage{xfrac}
\usepackage[unicode,hidelinks,bookmarks=false]{hyperref}
\newcommand{\ie}{i.e.\ }
\newcommand{\cf}{cf.\ }
\newcommand{\resp}{respectively\ }
\newcommand{\Subsection}{\subsection}
\providecommand{\nobreakdash}{\nobreak\hbox{-}}
\setlength{\emergencystretch}{2em}
\multlinegap0pt
\usepackage{bm}
\usepackage{bbm}
\usepackage{dsfont}
\usepackage{xspace}
\usepackage{cancel}
\usepackage{mathtools}
\usepackage{amsthm}
\makeatletter
\@ifundefined{theorem}{\newtheorem{theorem}{Theorem}}{}
\@ifundefined{lemma}{\newtheorem{lemma}[theorem]{Lemma}}{}
\makeatother
\title[The weak (1,1) boundedness of Fourier integral operators wit]{The weak (1,1) boundedness of Fourier integral operators with complex phases}
\author{Duv\'an Cardona, Michael Ruzhansky}
\date{Source: arXiv:2402.09054v5 (CC BY 4.0)}
\begin{document}
\maketitle

\noindent\emph{Source and licence.} The weak (1,1) boundedness of Fourier integral operators with complex phases. Version arXiv:2402.09054v5 (CC BY 4.0), distributed under the Creative Commons Attribution 4.0 International licence, as recorded in the arXiv licence metadata for this exact version and shown on the abstract page https://arxiv.org/abs/2402.09054v5. Full rights notice: licenses/arxiv-2402.09054-CC-BY-4.0.txt.\par\smallskip

\noindent\textit{Editorial scope.} Counted unit: the Lemma whose source label is \texttt{boundedness:A} in the article (source0.tex lines 980--982, source label \texttt{boundedness:A}), delivered together with the article's own complete original proof of it (source0.tex lines 983--1049, one \texttt{proof} environment of 67 lines). No mathematics is written, repaired or re-derived: statement and proof are the article's own text line for line. Required context retained uncounted: none. Every definition, display and label of those lines is retained unchanged. Omitted from this package and not redistributed: 1--979, 1050--1606. Nothing inside a retained excerpt is omitted or altered: every retained body is delivered exactly as the article's lines are written, so no omission receipt of any kind accompanies this package. Citations: the retained excerpts carry no citation command at all, so no bibliography entry of the article is transcribed and no reference is redistributed. Reference adaptation: none was needed and none was made; every reference inside the delivered unit resolves inside it, so no unresolved reference or citation is produced. Printed slips kept literal: no printed word, symbol, hypothesis or number of the counted statement or of its proof was altered. Layout and class: the article's own class is \texttt{amsart} with packages including amsmath, amsthm, amscd, amsfonts, amssymb, graphicx, color, mathrsfs, hyperref, xy, slashed, adjustbox, wrapfig; the wrapper is the lane's pinned \texttt{amsart} editorial wrapper at 10pt on A4, which restates the theorem-like environments the retained text needs and adds the article's own macro definitions that the retained text needs (none), with extra wrapper packages (bm, bbm, dsfont, xspace, cancel, mathtools). The delivered numbering is the unit's own. Assets: no retained span contains a figure environment, a graphics-inclusion command, a PDF-page inclusion, a table or a TikZ picture; no image file is copied into this package, the package declares no asset dependency and no image file is redistributed with it. Licence: version arXiv:2402.09054v5 of this article is distributed under the Creative Commons Attribution 4.0 International licence, as recorded in the arXiv licence metadata for this exact version and shown on the abstract page https://arxiv.org/abs/2402.09054v5. A full rights notice is delivered at licenses/arxiv-2402.09054-CC-BY-4.0.txt. Characters: every retained excerpt is pure ASCII, so the delivered engine, XeLaTeX with its default Latin Modern text font, reports no missing character.
\begin{lemma}\label{boundedness:A}
    $A$ extends to a bounded operator on $L^1(\mathbb{R}^n)$ provided that there exists $\varkappa_0>0,$ such that for all $(x,\theta)\in \overline{V}\times \overline{\{\theta: |\theta|\sim 1\}},$ the lower bound $$\textnormal{Im}(\Phi(x,\theta))>\varkappa_0,$$ holds.
\end{lemma}

\begin{proof}First, assume that $\widehat{f}(\xi)$ is smooth and compactly supported and that $f$ vanishes when $|\xi|\leq 1.$ Note that the support of $\eta_{k}(\xi')$ is contained in the region $|\xi'|\sim 2^ k.$
By applying summation by parts, because of the properties of the support of $f,$ boundary terms vanish and we can write
\begin{align*}
  &Af(x)\\
  &:= \sum_{k\gg 1}\smallint\smallint_{|\xi'|\sim 2^ k} e^{2\pi i\xi'\cdot \nabla_\xi\Phi_\tau(x,\omega)}\psi(x,\omega)e^{-\mu(x,\omega)\pi i/4} (1-\phi_{-\varepsilon k}(J(x,\omega)))\eta_{k}(\xi')e^{-(1+i\tau)\textnormal{Im}(\Phi(x,\xi'))}\\
&\times |J(x,\omega)|^{1/2}\widehat{f}(\xi')d\xi'd\omega\\
&=- \sum_{k\gg 1}\smallint\smallint_{|\xi'|\sim 2^ k} e^{2\pi i\xi'\cdot \nabla_\xi\Phi_\tau(x,\omega)}\psi(x,\omega)e^{-\mu(x,\omega)\pi i/4} (\phi_{-\varepsilon (k+1)}(J(x,\omega))-\phi_{-\varepsilon k}(J(x,\omega)))\\
&\times \phi_{k}(\xi')e^{-(1+i\tau)\textnormal{Im}(\Phi(x,\xi'))} 1_{\overline{V}}(x)|J(x,\omega)|^{1/2}\widehat{f}(\xi')d\xi'd\omega.
\end{align*}
In view of the triangle inequality, it suffices to prove  the estimate
\begin{equation}
    \Vert A_k f\Vert_{L^1(\mathbb{R}^n)}\lesssim 2^{-\varepsilon k}\Vert f\Vert_{L^1(\mathbb{R}^n)},
\end{equation} where 
\begin{align*}
    A_kf(x)&=- \smallint\smallint_{|\xi'|\sim 2^ k} e^{2\pi i\xi'\cdot \nabla_\xi\Phi_\tau(x,\omega)}\psi(x,\omega)e^{-\mu(x,\omega)\pi i/4}) (\phi_{-\varepsilon (k+1)}(J(x,\omega))-\phi_{-\varepsilon k}(J(x,\omega)))\\
&\times \phi_{k}(\xi')e^{-(1+i\tau)\textnormal{Im}(\Phi(x,\xi'))} 1_{\overline{V}}(x) |J(x,\omega)|^{1/2}\widehat{f}(\xi')d\xi'd\omega. 
\end{align*}
Let us denote by $$\phi_k^\bullet(\xi'):=\phi_k(\xi')\chi_{ \{\xi': |\xi'|\sim 2^k\} }(\xi').$$
Then,
\begin{align*}
    A_kf(x)&=- \smallint\smallint e^{2\pi i\xi'\cdot \nabla_\xi\Phi_\tau(x,\omega)}\psi(x,\omega)e^{-\mu(x,\omega)\pi i/4}) (\phi_{-\varepsilon (k+1)}(J(x,\omega))-\phi_{-\varepsilon k}(J(x,\omega)))\\
&\times \phi_{k}^\bullet(\xi')e^{-(1+i\tau)\textnormal{Im}(\Phi(x,\xi'))} 1_{\overline{V}}(x) |J(x,\omega)|^{1/2}\widehat{f}(\xi')d\xi'd\omega. 
\end{align*}
Using the Fourier inversion formula, we have that
\begin{align*}
   & \smallint e^{2\pi i\xi'\cdot \nabla_\xi\Phi_\tau(x,\omega)} \phi_{k}^\bullet(\xi') e^{-(1+i\tau)\textnormal{Im}(\Phi(x,\xi'))}1_{\overline{V}}(x) \widehat{f}(\xi')d\xi'\\
   &=1_{\overline{V}}(x)[e^{-(1+i\tau)\textnormal{Im}(\Phi(x,D))}\phi_{k}^\bullet(D)f](\nabla_\xi\Phi_\tau(x,\omega)).
\end{align*} Since $\textnormal{supp}[(\phi_{-\varepsilon (k+1)}(J(x,\omega))]$ is contained in the set where $J(x,\omega)\lesssim 2^{-\varepsilon k},$ the difference $\phi_{-\varepsilon (k+1)}(J(x,\omega))-\phi_{-\varepsilon k}(J(x,\omega)),$ is supported on the region $J(x,\omega)\sim 2^{-k\varepsilon}.$ Since $|\phi_k(t)|\lesssim \Vert\phi\Vert_{L^{\infty}},$ $\forall t>0,$ we have that $$|(\phi_{-\varepsilon (k+1)}(J(x,\omega))-\phi_{-\varepsilon k}(J(x,\omega))||J(x,\omega)|^{1/2}\lesssim 2^{-\varepsilon k/2},$$
and consequently
\begin{align*}
    \Vert A_k f\Vert_{L^1(\mathbb{R}^n)}  &=\smallint|A_k f(x)|dx\\
    &\lesssim\smallint\smallint |1_{\overline{V}}(x)[e^{-(1+i\tau)\textnormal{Im}(\Phi(x,D))}\phi_{k}^\bullet (D)f](\nabla_\xi\Phi_\tau(x,\omega))|2^{-\varepsilon k/2}|\phi(x,\omega)|dx d\omega.
\end{align*}    Note that the operator 
\begin{equation*}
   \Omega(x,D)= 1_{\overline{V}}(x)e^{-(1+i\tau)\textnormal{Im}(\Phi(x,D))}\phi_{k}^\bullet(D)
\end{equation*}is uniformly bounded in $k$ on $L^1(\mathbb{R}^n).$ Indeed, its kernel $\Omega(x,D)\delta(y),$ given by  
$$ \Omega(x,D)\delta(y)=\smallint e^{2\pi i(x-y)\xi}1_{\overline{V}}(x)\phi_{k}^\bullet(\xi)e^{-(1+i\tau)\textnormal{Im}(\Phi(x,\xi))}d\xi  $$
belongs to $L^1(\mathbb{R}_x^n)$ uniformly in $y\in \mathbb{R}^n.$
Since  the phase function $\Phi$ is of strictly positive type, that is, $\textnormal{Im}(\phi)> 0,$ and $|\phi_{k}(\xi)|\lesssim 1,$ note that
\begin{align*}
    \smallint|\smallint e^{2\pi i(x-y)\xi}1_{\overline{V}}(x) &\phi_{k}^\bullet(\xi)e^{-(1+i\tau)\textnormal{Im}(\Phi(x,\xi))}d\xi|dx \\
   = \smallint|\smallint_{|\xi|\sim 2^k} e^{2\pi i(x-y)\xi}1_{\overline{V}}(x) &\phi_{k}(\xi)e^{-(1+i\tau)\textnormal{Im}(\Phi(x,\xi))}d\xi|dx \\
    &\leq \smallint   1_{\overline{V}}(x)\sup_{|\xi'|\sim 2^{k}}e^{-\textnormal{Im}(\Phi(x,\xi'))} dx\smallint_{|\xi|\sim 2^{k}}|\phi_{k}(\xi)| d\xi\\
    &\lesssim |\overline{V}|\sup_{x\in \overline{V}} \sup_{|\xi'|\sim 2^{k}}e^{-\textnormal{Im}(\Phi(x,\xi'))}\smallint_{|\xi|\sim 2^{k}}d\xi\\
    &\lesssim |\overline{V}|\sup_{x\in \overline{V}} \sup_{|\xi'|\sim 2^{k}}e^{-2^{k}\textnormal{Im}(\Phi(x,2^{-k}\xi'))}2^{kn}\\
    & \lesssim |\overline{V}|\sup_{x\in \overline{V}} \sup_{|\theta|\sim 1}e^{-2^{k}\textnormal{Im}(\Phi(x,\theta))}2^{kn}<\infty.
\end{align*}Indeed, since $\textnormal{Im}(\phi)> 0,$ when $|\xi|\neq 0,$ the compactness of the closed annulus $\overline{\{\theta: |\theta|\sim 1\}},$ and of $\overline{V},$ imply that there exists $\varkappa_0>0$ such that for all $(x,\omega)\in \overline{V}\times \overline{\{\theta: |\theta|\sim 1\}},$ one has the lower bound $$\textnormal{Im}(\Phi(x,\theta))>\varkappa_0.$$ Therefore, 
$$\mathscr{K}:=\sup_{k>0}\sup_{x\in \overline{V}} \sup_{|\theta|\sim 1}e^{-2^{k}\textnormal{Im}(\Phi(x,\theta))}2^{kn}\leq \sup_{k>0} e^{-2^{k}\varkappa_0}2^{kn}<\infty.$$
Thus, we have proved that  $\sup_{y\in \mathbb{R}^n}\Vert  \Omega(x,D)\delta(y) \Vert_{L^1(\mathbb{R}^n_x)}<\infty,$ and Schur's lemma gives the boundedness of $\Omega(x,D)$ on $L^1(\mathbb{R}^n).$ Note also, that if $z=F(x)$ is an smooth change of coordinates on $\overline{V}$, the operator $$\Omega(F^{-1}(z),D)=1_{\overline{V}}(F^{-1}(z))\phi_{k}^\bullet(D)e^{-(1+i\tau)\textnormal{Im}(\Phi(F^{-1}(z),D))}$$ is also bounded on $L^1.$ Moreover, since $1_{\overline{V}}\circ F^{-1}=1_{F(\overline{V})},$ by following the previous argument, we have that
\begin{align*}
   \sup_{y}  \Vert \Omega(F^{-1}(z),D)\delta(y)\Vert_{L^1(\mathbb{R}^n_z)}\lesssim  |F(\overline{V})|\sup_{z\in F(\overline{V})} \sup_{|\theta|\sim 1}e^{-2^{k}\textnormal{Im}(\Phi(F^{-1}(z),\theta))}2^{kn}<\infty.
\end{align*}Now, since for any $\omega,$ $z=F_\omega(x)=\nabla_\xi\Phi_\tau(x,\omega)$ is a diffeomorphism on $\overline{V},$ taking $R>0$ large enough such that $\{\omega: (x,\omega)\in \textnormal{supp}(\phi(x,\omega))\}\subset\{|\omega|\leq R\},$ we have that
\begin{align*}
    &\Vert A_k f\Vert_{L^1(\mathbb{R}^n)}  \\
    &\lesssim\smallint_{|\omega|\leq R}\smallint |1_{\overline{V}}(x)[\phi_{k}^\bullet(D)e^{-(1+i\tau)\textnormal{Im}(\Phi(x,D))}f](\nabla_\xi\Phi_\tau(x,\omega))|2^{-\varepsilon k/2}|\phi(x,\omega)|dx d\omega\\
    &\lesssim \Vert\phi(x,\omega)\Vert_{L^{\infty}} \smallint_{|\omega|\leq R}\smallint |\Omega(x,D)f(F_{\omega}(x))|2^{-\varepsilon k/2}dx d\omega\\
    &= \Vert\phi(x,\omega)\Vert_{L^{\infty}} \smallint_{|\omega|\leq R}\smallint |\Omega(F^{-1}_\omega(z),D)f(z)|2^{-\varepsilon k/2}|\det(DF_{\omega}^{-1}(z))|dz d\omega\\
    &\lesssim 2^{-\varepsilon k/2} \smallint_{|\omega|\leq R}\Vert\Omega(F^{-1}_\omega(z),D)f\Vert_{L^1(\mathbb{R}^n_z)}d\omega\\
    &\lesssim 2^{-\varepsilon k/2} \sup_{|\omega'|\leq R}\Vert\Omega(F^{-1}_{\omega'}(z),D)f\Vert_{L^1(\mathbb{R}^n_z)} \smallint_{|\omega|\leq R}d\omega\\
    & \lesssim 2^{-\varepsilon k/2} \sup_{|\omega'|\leq R} |F_{\omega'}(\overline{V})|\sup_{z\in F_{\omega'}(\overline{V})} \sup_{|\theta|\sim 1}e^{-2^{k}\textnormal{Im}(\Phi(F_{\omega'}^{-1}(z),\theta))}2^{kn}\Vert f\Vert_{L^1(\mathbb{R}^n)}\\
    &=2^{-\varepsilon k/2} \sup_{|\omega'|\leq R} |F_{\omega'}(\overline{V})|\sup_{x\in \overline{V}} \sup_{|\theta|\sim 1}e^{-2^{k}\textnormal{Im}(\Phi(x,\theta))}2^{kn}\Vert f\Vert_{L^1(\mathbb{R}^n)}\\
    &\lesssim_{\mathscr{K}} 2^{-\varepsilon k/2} \sup_{|\omega'|\leq R} |F_{\omega'}(\overline{V})|\Vert f\Vert_{L^1(\mathbb{R}^n)}.
\end{align*} Since $F_{\omega}$ is $C^1$ in $\omega,$ and $\overline{V}$ is compact, $ \sup_{|\omega'|\leq R} |F_{\omega'}(\overline{V})|<\infty.$ So we have the estimate
\begin{equation}
    \Vert A_k f\Vert_{L^1(\mathbb{R}^n)}\lesssim 2^{-k\varepsilon/2}\Vert f\Vert_{L^1(\mathbb{R}^n)}
\end{equation}as desired proving the $L^1$-boundedness of $A.$ The proof of Lemma \ref{boundedness:A} is complete.
\end{proof} 

\end{document}
